\section{Nonnormal generators and the limits of the comparison}\label{sec:examples}
We first record an automatic-continuity fact used to exclude every algebraic full-domain system in the first separation.

\begin{lemma}\label{lem:automatic}
Let $x\in\M$ be a bounded normal operator and $E:\M\Omega\to H$ a complex-linear map satisfying $[E,x]=P_\Omega$.
Then $E$ is bounded from the operator norm on $\M$ to $H$.
No adjoint or closability hypothesis on $E$ is required.
\end{lemma}
\begin{proof}
The spectral-atom argument in Lemma~\ref{lem:commutation} gives a diffuse law.
On the Banach space $H\oplus\M$, put
\[
 R:=\begin{pmatrix}N&P_\Omega|_\M\\0&T\end{pmatrix},\qquad
 N:=\lambda(x)|_H,\quad T(a):=xa,\quad J(a):=(Ea,a).
\]
The dual identity is $RJ=JT$.
We verify the automatic-intertwining hypotheses of \cite[Proposition~3.7 and Theorem~4.1]{Lau88}.
For a smooth function $\phi$ on a neighborhood of $\sigma(x)$, define
\[
 \Phi(\phi):=
 \begin{pmatrix}\phi(N)&\Gamma_\phi|_\M\\0&L_{\phi(x)}\end{pmatrix},\qquad
 \Gamma_\phi:=\iint\frac{\phi(z)-\phi(w)}{z-w}\,dE_x(z)P_\Omega\,dE_x(w).
\]
The integral is Hilbert--Schmidt, with norm bounded by the Lipschitz norm on a fixed compact neighborhood.
Its product-null diagonal causes no ambiguity.
The divided-difference product rule gives a multiplicative calculus with $\Phi(1)=1$ and $\Phi(z)=R$.
Given a finite open spectral cover, a smooth partition of unity supplies commuting splitting maps whose closed ranges are $R$-invariant and have spectra inside the corresponding members of the cover.
For the spectral containment, multiply a cutoff by a smooth local reciprocal $(z-\lambda)^{-1}$ outside its compact support.
This is the super-decomposability condition used in the cited proposition.
The operator $T$ is decomposable, by its left spectral projections on $\M$.

There is no nonzero algebraically divisible subspace for $N$.
Indeed, a vector $v$ in such a subspace belongs to $\Ran(N-\lambda)^2$ for every $\lambda$, so its spectral measure satisfies $I_v(\lambda):=\int|z-\lambda|^{-4}\dd\nu_v(z)<\infty$ everywhere.
If $\nu_v\ne0$, some set $A_n:=\{I_v\le n\}$ has positive $\nu_v$ measure.
For centers in $A_n$, $\nu_v(B(\lambda,r))\le nr^4$.
A grid covers bounded $A_n$ by $O(r^{-2})$ disks of radius $O(r)$ centered in $A_n$, forcing $\nu_v(A_n)\le cnr^2$, a contradiction.
The inclusion $\M\hookrightarrow H$ rules out divisible subspaces for $T$ as well.
Projecting a putative $R$-divisible subspace to its second coordinate and then its first rules it out for $R$.
By \cite[Proposition~3.7]{Lau88}, its algebraic spectral subspaces for closed sets are closed.
Both $R$ and $T$ have no eigenvalues, so there are no critical eigenvalues for this pair.
The intertwining theorem \cite[Theorem~4.1]{Lau88} therefore makes $J$ continuous.
Its first component is $E$.
\end{proof}

\subsection{Conjugate variables without any algebraic full-domain system}
For a bounded single generator $T$ define
\[
 \beta(T):=\inf\{\|E\|_{\M\to H}:[E,T]=P_\Omega\text{ on }\M\Omega\},
\]
with infinite value if there is no bounded witness.
The value $E\Omega$ is unrestricted.
Let $G:=\langle u,s\mid s^2=1\rangle=\mathbb Z*\mathbb Z_2$ and work in its group von Neumann algebra with the identity trace.
For $0<\varepsilon\le1/2$, put
\[
 p:=(1+s)/2,\quad q:=1-p,\quad h:=\varepsilon p+q,\quad T:=uh.
\]
Since $T^*T=\varepsilon^2p+q$, its represented star algebra contains $p,s,h^{-1}$ and $u=Th^{-1}$ and is exactly $\C G$.
In particular $T$ generates the algebra and has norm one.
Writing $t:=usu^{-1}$, prescribe
\begin{equation}\label{eq:group-derivative}
 \delta u=1\otimes q+\varepsilon^{-1}t\otimes p,
 \qquad \delta p=-\frac2{1-\varepsilon}qu^{-1}\otimes p,
 \qquad \delta(u^{-1})=-u^{-1}(\delta u)u^{-1}.
\end{equation}
The prescriptions respect $s^2=1$ and both inverse relations, so they define an outer derivation on the reduced group words.
They give $\delta T=1\otimes1$.
Its actual vacuum is
\begin{equation}\label{eq:group-vacuum}
 \xi=\tfrac12u-\left(\tfrac12+\frac1{1-\varepsilon}\right)us,
 \qquad \|\xi\|_2^2\le13/2.
\end{equation}
Here is a check on all words.
For a marked positive $u$ cut, the term $t\otimes p$ would require a prefix ending in $u^{-1}$, which cannot precede $u$ in a reduced word.
The term $1\otimes q$ contributes $1/2$ at $u$ and $-1/2$ at $us$.
For a negative $u^{-1}$ cut, a nonzero double trace would require a suffix starting in $u$, which is likewise impossible.
For an $s$ cut, expand $\delta s=-(u^{-1}-su^{-1})\otimes(1+s)/(1-\varepsilon)$.
All cuts except the prefix $u$ and empty suffix would force adjacent $ss$; that cut contributes $-1/(1-\varepsilon)$ at $us$.
This exhausts every marked occurrence in every reduced word and proves the actual pairing in \eqref{eq:group-vacuum}.

For any bounded witness $E$, let $f:=upu^{-1}$.
Then $\|Tp\|=\|fT\|=\varepsilon$ and $\tau(f)=\tau(p)=1/2$.
Evaluating the commutator at $p$ and multiplying by $f$ gives
\begin{equation}\label{eq:dual-cost}
 \frac1{2\sqrt2}=\|\tau(p)f\|_2
 \le\|fE(Tp)\|_2+\|fTE(p)\|_2
 \le2\varepsilon\|E\|_{\M\to H}.
\end{equation}
Thus $\beta(T)\ge1/(4\sqrt2\varepsilon)$, uniformly over all witnesses.

We now assemble these increasing costs into a single actual law.
Choose a sequence $T_n$ of the above examples for which
\[
 \beta(T_n)/(1+\|\xi_n\|_2)\ge4^n,
\]
and let $K_n:=\delta_n(T_n^*)$.
Choose positive weights $t_n$ so small that
\[
 t_n\le2^{-3n-4},\qquad t_n\|K_n\|_2\le2^{-n},\qquad
 r_n:=2^nt_n^{3/2}(1+\|\xi_n\|_2)\le2^{-n-4}.
\]
Put $t_0:=1-\sum_{n\ge1}t_n>0$ and use a fixed norm-one semicircular input at index zero.
In the product algebra $\M:=\prod_{n\ge0}\M_n$ with trace $\tau:=\sum_nt_n\tau_n$, set
\[
 X_n:=2^{-n}1_n+r_nT_n\ (n\ge1),\qquad
 X_0:=2\cdot1_0+\tfrac1{16}T_0,\qquad X:=(X_n)_n.
\]
The disjoint spectral disks recover each component unit by a contour integral.
Uniform Neumann resolvent bounds outside their closed union justify this calculus in the product algebra.
Thus $X$ is bounded and generates $\M$.

On an evaluated star polynomial define the diagonal coarse block of $\delta p(X)$ to be $r_n^{-1}\delta_n(p(X_n))$; use $r_0=1/16$ at zero.
The diagonal generator values are $1_n\otimes1_n$ at $X$ and $K_n$ at $X^*$, and they are square summable with weights $t_n^2$.
The diagonal vacuum has component $t_n\xi_n/r_n$ and finite squared norm, since
\[
 \sum_{n\ge1}\frac{t_n^3}{r_n^2}\|\xi_n\|_2^2\le\sum_{n\ge1}4^{-n}.
\]
On each mixed coarse block choose $k_{ab}$ solving
\[
 X_ak_{ab}-k_{ab}X_b=1_a\otimes1_b\qquad(a\ne b).
\]
The separated centers and $r_a+r_b\le|\lambda_a-\lambda_b|/2$ give a Neumann solution with
\[
 \|k_{ab}\|_2\le\frac{2\sqrt{t_at_b}}{|\lambda_a-\lambda_b|},
 \qquad \sum_{a\ne b}\|k_{ab}\|_2^2<\infty,
\]
where $\lambda_n=2^{-n}$ for $n\ge1$ and $\lambda_0=2$.
Let $k$ be this direct sum and prescribe the mixed derivative by $p\mapsto p(X)k-kp(X)$.
It has actual vacuum norm at most $2\|k\|_2$.
Combining the diagonal and mixed blocks gives a well-defined represented derivation with $\delta X=1\otimes1$ and actual $L^2$ vacuum.
Consequently $\mathsf A(X)$ holds.

No bounded witness at $X$ exists.
Indeed, compress any $E:\M\to H$ to input and output corner $n$, measured in normalized norms, and call it $B_n$.
The corner relation gives $[B_n,T_n]=(t_n/r_n)P_{1_n}$.
Independent signs on finitely many corner inputs of norm at most one give
\[
 \|E\|^2\ge\sum_nt_n\|B_n\|^2
 \ge\sum_{n\ge1}\frac{t_n^3}{r_n^2}\beta(T_n)^2
 =\sum_{n\ge1}4^{-n}\frac{\beta(T_n)^2}{(1+\|\xi_n\|_2)^2}=\infty.
\]
The sign average cancels the cross terms and therefore also controls witnesses with arbitrary off-diagonal blocks.

Freely adjoin a standard semicircular operator $s_0$ and consider the pair $(X,s_0)$.
The old derivation extends by zero on $s_0$, preserving its actual vacuum by the centered reduced-word calculation.
The semicircular difference quotient extends by zero on the old algebra, with actual vacuum $s_0$.
Thus the pair satisfies $\mathsf A$.
If it had an algebraic system, change signs to maps $E_X,E_s$ with the operator-first convention.
Both $E_s$ and $E_s+E_X$ have commutator $P_\Omega$ with the normal operator $s_0$.
Lemma~\ref{lem:automatic} makes both bounded for the operator norm, hence $E_X$ bounded.
Restriction to the old algebra and conditional expectation onto its $L^2$ space would give a bounded witness at $X$, a contradiction.
This excludes all algebraic full-domain systems and proves the first separation in Theorem~\ref{thm:general}.

\subsection{A bounded strong dual without a starred derivation}
Let
\[
 d\mu(s):=\pi^{-1}\sqrt{4-s^2}\,1_{(0,2)}(s)\,ds,
 \qquad\M:=M_2(L^\infty((0,2),\mu)),
 \qquad\tau(a):=\tfrac12\int\Tr_2(a(s))\dd\mu(s).
\]
Put
\[
 J_0:=\begin{pmatrix}1&0\\0&-1\end{pmatrix},\quad
 h:=\begin{pmatrix}1&-1/2\\0&1\end{pmatrix},\quad
 K:=hJ_0h^{-1}=\begin{pmatrix}1&1\\0&-1\end{pmatrix},
 \qquad Y:=sJ_0,\quad X:=sK.
\]
The bounded nonnormal operator $X$ generates $\M$.
Indeed, $X^2=s^2I_2$ recovers the scalar center and bounded truncations of $s^{-1}X$ recover $K$.
With $q:=(K+I_2)/2$ one has $e_{11}=4qq^*/5$, $e_{12}=2(q-e_{11})$, and hence all constant matrix units.

The operator $Y$ has the standard semicircle law, by integrating $f(s)$ and $f(-s)$ with weight $1/2$.
On $L^2(W^*(Y))$ its polynomial basis identifies $Y$ with $\ell+\ell^*$ and gives $[\ell^*,Y]=P_\Omega$ \cite{Miy22}.
This subspace reduces $Y$; extend $\ell^*$ by zero on its orthogonal complement to obtain $E_Y\in\mathcal B(H)$ with the same commutator on all $H$.
The bounded invertible map $W(a\Omega):=hah^{-1}\Omega$ intertwines $Y$ and $X$.
It fixes $\Omega$ and preserves the vacuum functional by traciality, so $WP_\Omega W^{-1}=P_\Omega$.
Consequently $D:=-WE_YW^{-1}$ is bounded on all $H$ and satisfies $[X,D]=P_\Omega$.
It witnesses $\mathsf S(X)$.

Suppose, on the other hand, that an outer derivation on $\C\langle X,X^*\rangle$ with $\delta X=1\otimes1$ existed, and put $\beta:=\delta X^*$.
Differentiate the represented relation $[X^2,X^*]=0$ to obtain
\[
 (X^2\otimes1-1\otimes X^2)\beta
 =X^*\delta(X^2)-\delta(X^2)X^*,\qquad
 \delta(X^2)=X\otimes1+1\otimes X.
\]
In the two scalar coordinates $s,t$, this says
\[
 (s^2-t^2)\beta(s,t)
 =s^2K^*K\otimes I_2+stK^*\otimes K-stK\otimes K^*-t^2I_2\otimes KK^*.
\]
Extracting the $e_{21}\otimes e_{11}$ coefficient gives $(s^2-t^2)b(s,t)=s(s+t)$, and hence $b(s,t)=s/(s-t)$ almost everywhere.
This function is not in $L^2(\mu\otimes\mu)$, since the density is bounded below on $[1/2,1]^2$.
Coefficient extraction is bounded on the finite-dimensional matrix tensor factor, contradicting $\beta\in H\bar\otimes H$.
This excludes every normalized starred Hilbert-valued derivation before any vacuum condition is imposed.
Together with Proposition~\ref{prop:bare-maximal} and \cite[Proposition~3.5]{ACSY24}, it completes Theorem~\ref{thm:general}.
